\chapter{Step 7 Formulation Strengthening Study}
\label{ch:step7}

This appendix records an isolated formulation-strengthening experiment motivated by the valid-inequality and reformulation agenda identified in Step~7 of the research roadmap. The baseline model and the corrected Constraint~(13) presented in the main thesis are retained unchanged. The purpose of this appendix is to evaluate an additional trigger-to-assignment linking inequality without incorporating it into the baseline formulation.

\section{Strengthened trigger linking}

For a maintenance-eligible flight $i$, aircraft $j$, check type $c$, and sparse set of feasible trigger days $D(i,c)$, the strengthened variant imposes

\begin{equation}
    \sum_{d\in D(i,c)} z_{ijdc} \leq x_{ij},
    \qquad \forall i,j,c.
    \label{eq:step7_strong_link}
\end{equation}

The inequality states that no maintenance trigger associated with flight $i$ and aircraft $j$ can be active unless that flight is assigned to that aircraft. Unlike the baseline C9 implementation for calendar-day checks, which links only the first candidate trigger day, equation~\eqref{eq:step7_strong_link} links every sparse candidate trigger day. The constraint is implemented as an opt-in variant through the
\texttt{use\_strong\_maint\_link} flag in \texttt{src/model.py}; its default value is \texttt{False}.

The strengthening is logically redundant for integer solutions whenever the remaining maintenance constraints already prevent an unassigned flight from activating a deferred trigger. It can nevertheless improve the LP relaxation by removing fractional trigger mass that is not supported by the corresponding flight-assignment variable. The expected benefit is therefore a stronger bound and, potentially, fewer branch-and-bound nodes rather than a change in the intended integer feasible set.

\section{Implementation and test protocol}

The baseline and strengthened variants use identical flight data, aircraft data, maintenance parameters, objective coefficients, solver settings, and time limits. The comparison harness in \texttt{experiments/compare\_strong\_maint\_link.py} records model size, C9 row count, sparse $z$-domain size, build time, LP objective, solver status, and solve time in a separate Step~7 result file.

Two structural regression tests were added. The first creates a controlled two-day type-C maintenance window and verifies that the strengthened constraint contains both candidate $z$ variables, whereas the baseline links only the first candidate day. The second verifies that omitting the new flag is exactly equivalent to explicitly selecting the baseline formulation. These tests validate the formulation distinction without requiring a solver to infer the intended logic from an optimal solution.

\section{Initial comparison}

The first comparison was executed on the existing smoke instance
\texttt{DataCplex\_density=0.5\_p=4\_h=7\_test\_0}. Both variants contain 984 variables, 1,964 constraints, 224 C9 rows, and 616 sparse $z$ variables. The identical counts are expected for this instance because its calendar-day maintenance windows contain only one candidate trigger day. The synthetic multi-day regression, however, produces the same number of C9 rows in both variants while changing each relevant C9 expression from a single-day link to a sum over all feasible trigger days. The improvement is therefore structural and is activated only when deferred multi-day trigger windows exist.

\begin{table}[htbp]
\centering
\caption{Initial structural comparison of the Step~7 linking variant.}
\label{tab:step7_structure}
\begin{tabular}{lrrrr}
\toprule
Variant & Variables & Constraints & C9 rows & $z$ variables \\
\midrule
Baseline & 984 & 1,964 & 224 & 616 \\
Strong linking & 984 & 1,964 & 224 & 616 \\
\bottomrule
\end{tabular}
\end{table}

These initial results do not establish a runtime improvement. They show instead that the strengthening preserves model size on the current smoke instance and that its discriminating effect is concentrated in the LP expression associated with multi-day deferred triggers. A larger matched experiment containing long horizons and multi-day maintenance windows is required before making a computational superiority claim. Accordingly, this appendix reports the implementation as a separate strengthening claim and does not alter the conclusions of the main formulation chapters.

\section{Controlled LP comparison}

To obtain a solver-backed comparison below the installed CPLEX Community
Edition size limit, a minimal two-flight, one-aircraft fixture was generated
with a two-day type-C maintenance duration. Both formulations were relaxed by
changing binary variables to continuous variables in $[0,1]$ and were solved
with identical CPLEX settings. The two LPs reached optimality and returned the
same objective value, while the strengthened variant required less wall-clock
time on this particular run.

\begin{table}[htbp]
\centering
\caption{Controlled LP comparison for the Step~7 strengthening.}
\label{tab:step7_lp_control}
\begin{tabular}{lrrrrr}
\toprule
Variant & Variables & Constraints & C9 rows & LP objective & Runtime (s) \\
\midrule
Baseline & 85 & 211 & 4 & 20.0 & 0.0798 \\
Strong linking & 85 & 211 & 4 & 20.0 & 0.0459 \\
\bottomrule
\end{tabular}
\end{table}

The equal LP objective values are expected for this deliberately small
fixture: the trigger-linking inequality removes unsupported fractional trigger
patterns but does not alter the optimum of this instance. The runtime
difference is descriptive only and must not be interpreted as a general speed
claim. A statistically meaningful evaluation requires repeated matched
instances with multi-day maintenance windows, larger horizons, and both LP and
integer solves. The corresponding machine-readable result is stored in
\texttt{results/tables/step7\_strong\_link\_lp\_tiny\_control.csv}.
